% ============================================================
%  chapters/05-conclusion/conclusion.tex
%  Conclusion
% ============================================================

\chapter*{Conclusion}
\addcontentsline{toc}{chapter}{Conclusion}
\markboth{Conclusion}{Conclusion}

% ── Recall of context ────────────────────────────────────────
This report was carried out within \CompanyName{} during the period
\InternshipPeriod{} as part of the final-year engineering programme
at \FacultyName{}. The work addressed a concrete industrial challenge
facing organisations that operate or plan to operate AI infrastructure:
the near-total dependency of modern large language model pipelines on
NVIDIA's proprietary CUDA ecosystem, and the technical feasibility of
transitioning those pipelines to Huawei's Ascend CANN ecosystem as a
sovereign, energy-efficient alternative.

% ── Summary of contributions ─────────────────────────────────
\section*{Summary of Contributions}
\addcontentsline{toc}{section}{Summary of Contributions}

The work unfolded across four structured chapters, each building
on the previous to form a coherent end-to-end study.

\textbf{Chapter~1} established the general context of the project:
the organisational context of \CompanyName{} and its position within
the Huawei group, the technical context of the AI hardware landscape,
and the strategic motivations — supply chain resilience, energy
efficiency, and sovereignty — that make GPU-to-NPU migration a
relevant and pressing challenge for enterprises operating in the
MENA region.

\textbf{Chapter~2} provided the technological foundations necessary
to understand the migration problem at a systems level. The chapter
covered GPU computing foundations — the CUDA execution model,
the SIMT architecture, and the CUDA software ecosystem — followed
by NPU architecture principles and the DaVinci compute model, and
concluded with a detailed analysis of the Huawei Ascend stack:
the DaVinci core architecture, the CANN runtime and compiler,
the MindSpore framework, and the complete AI software and
infrastructure landscape spanning both the NVIDIA and Huawei
ecosystems.

\textbf{Chapter~3} conducted a technical reverse-engineering study
of the Nemotron 3 Nano pipeline as a concrete, production-grade
case study of an NVIDIA-native AI lifecycle. The chapter
reconstructed the complete training pipeline from the repository
— data preparation, pretraining, supervised fine-tuning,
reinforcement learning, evaluation, and inference — and produced
a component-level technical stack analysis identifying every
CUDA dependency, its position in the pipeline, and its functional
role. This analysis established the precise scope of the migration
challenge: a vertically integrated CUDA stack with two
architecture-specific custom operators, a complex distributed
training framework, and a tight coupling between the training
and inference stages through shared model formats and serving
infrastructure.

\textbf{Chapter~4} translated the findings of Chapter~3 into a
structured, actionable migration plan. The plan addressed the
complete scope of the migration — from hardware provisioning
through operator implementation, framework porting, and inference
deployment — organised into four sequential phases with explicit
acceptance gates. Two implementation tracks were defined to
separate correctness validation from performance optimisation.
A component-level change impact assessment identified which
files require no changes, which require targeted modifications,
and which represent genuine engineering work. A risk register
and validation framework completed the plan, providing the
operational structure needed to execute the migration as a
managed infrastructure project.

% ── Key findings ─────────────────────────────────────────────
\section*{Key Findings}
\addcontentsline{toc}{section}{Key Findings}

Several conclusions of broader relevance emerged from this study.

\textbf{The migration complexity is not uniform.} A common
assumption about GPU-to-NPU migration is that it requires
rewriting an AI pipeline from scratch. This study demonstrates
that the reality is more nuanced: a substantial portion of a
production AI pipeline — the data preparation infrastructure,
the experiment tracking layer, the orchestration system,
and the distributed communication primitives — is entirely
hardware-agnostic and requires no changes whatsoever. The
migration complexity concentrates at a small number of specific
layers: the compute runtime, the training framework's
device integration, and, critically, any custom hardware
kernels specific to the model architecture.

\textbf{Architecture-specific operators are the critical path.}
For Nemotron 3 Nano, the two custom CUDA operators required
by the Mamba SSM and MoE components represent the highest-risk,
highest-effort items in the entire migration — more so than
framework porting or infrastructure reconfiguration. This
finding generalises: any model that relies on custom hardware
kernels beyond the standard set of neural network primitives
will face a per-operator implementation bottleneck on any
non-CUDA platform, regardless of how mature the broader
ecosystem is.

\textbf{The PyTorch-NPU compatibility layer substantially reduces
migration friction.} Huawei's \texttt{torch\_npu} package, which
registers Ascend NPUs as a native PyTorch device type, means that
the majority of PyTorch-based training code can run on Ascend
with targeted rather than wholesale changes. This significantly
lowers the initial barrier to entry and enables a correctness
validation baseline to be established quickly before committing
to a deeper MindSpore-native reimplementation.

\textbf{The energy efficiency advantage is material at scale.}
The Ascend 910C's performance-per-watt advantage over the H100
— 2.00 versus 1.41 TFLOPS/W — translates into a 42\% reduction
in energy consumption for equivalent compute work. At the scale
of a data centre running thousands of training hours per month,
this efficiency gain represents a substantial and persistent
reduction in operating expenditure that compounds over the
infrastructure lifetime, partially or fully offsetting the
one-time migration investment.

\textbf{Ecosystem maturity remains the primary adoption barrier.}
Despite the hardware performance and efficiency advantages of
Ascend 910C, the CANN ecosystem's maturity gap relative to CUDA
— in terms of operator coverage, developer tooling, community
resources, and production-validated framework integrations —
represents the largest practical barrier to adoption. Projects
that rely exclusively on standard transformer architectures are
substantially better positioned for migration than those that
use novel architectural primitives, as was the case with
Nemotron 3 Nano's Mamba SSM layers.

% ── Limitations ──────────────────────────────────────────────
\section*{Limitations}
\addcontentsline{toc}{section}{Limitations}

Several limitations of this study should be acknowledged.

The migration plan presented in Chapter~4 is a design and
planning artefact: the full implementation, numerical validation,
and performance benchmarking of the migrated pipeline on physical
Ascend 910C hardware falls outside the scope of this report.
The performance estimates in Chapter~4 are derived from hardware
specifications and published benchmarks rather than from
empirical measurements on the target workload. Actual throughput
and latency figures may differ, particularly for the
architecture-specific Mamba SSM and MoE operations whose
performance under TBE implementation is inherently speculative
at this stage.

Furthermore, the study focuses on a single model architecture
— Nemotron 3 Nano — as the case study. While the migration
methodology is designed to be generalisable, the specific
operator coverage gaps, the parallelism configuration
adjustments, and the checkpoint conversion requirements are
architecture-dependent and would need to be reassessed for
a different model family.

Finally, the vLLM-Ascend fork used as the generation backend
for the RL stage is a community-maintained project whose
production readiness for MoE architectures and long-context
configurations was not empirically validated within this work.
Its suitability for the full Nemotron 3 Nano production
deployment remains to be confirmed through practical testing.

% ── Future work ──────────────────────────────────────────────
\section*{Perspectives and Future Work}
\addcontentsline{toc}{section}{Perspectives and Future Work}

Several directions for future work emerge naturally from this
study.

\textbf{Empirical validation and benchmarking.} The most
immediate next step is the execution of the migration plan:
implementing the custom TBE operators, patching the training
scripts, and running the validation gate sequence on physical
Ascend 910C hardware. This would produce empirical training
throughput, inference latency, and benchmark accuracy figures
that could confirm or revise the estimates presented in
Chapter~4, and would constitute a complete proof of concept
for the migration approach.

\textbf{TBE operator optimisation.} The TBE implementations
of the Mamba SSM selective scan and MoE token dispatch operators
represent significant optimisation opportunities. Initial
implementations using straightforward approaches are likely
to leave substantial performance on the table relative to what
is achievable with careful tiling, pipelining, and DaVinci
Cube/Vector unit co-scheduling. A dedicated operator
engineering effort focusing on these two primitives could
significantly close the throughput gap between Ascend and H100.

\textbf{Generalisation to a migration intelligence platform.}
The methodology developed in this report — systematic
dependency analysis, operator coverage mapping, phased migration
planning, and gate-based validation — could be formalised into
a reusable migration intelligence tool. Such a tool would
accept a model architecture and a target hardware platform
as inputs, automatically identify coverage gaps, generate a
migration effort estimate, and suggest an optimal migration
track. This was the original framing of the project and
represents a high-value direction for continued development
at \CompanyName{}.

\textbf{Broader hardware target coverage.} The methodology
applied in this report to the CUDA-to-CANN migration is
applicable to other hardware transitions: AMD ROCm, Intel
oneAPI, and custom ASIC accelerators. Extending the analysis
to cover these platforms would produce a comprehensive
multi-vendor migration framework relevant to any organisation
seeking to diversify its AI hardware dependency beyond the
current NVIDIA-dominated market.

\textbf{Agentic and long-context deployment on Ascend.}
Nemotron 3 Nano's 1M-token context capability and its
agentic tool-calling features represent deployment scenarios
that place particularly demanding requirements on the
inference stack. Validating the complete MindIE inference
pipeline at 1M-token context and under multi-turn agentic
workloads — with accurate measurement of latency, throughput,
and KV-cache efficiency — would provide a definitive assessment
of Ascend's readiness for the most demanding LLM deployment
patterns.

\vspace{1.5cm}

\begin{flushright}
  \textit{\AuthorFullName}\\
  \ReportDate{}~—~\ReportLocation
\end{flushright}
